\chapter{시간: 언제 무엇을 하는가}\label{ch:relay-timing}

\ref{ch:sync}장의 동기식 회로에는 시간 규칙이 하나뿐이었다. ``한 클럭 안에 다음 레지스터에 도착하라.''
설계 도구가 모든 경로의 지연을 계산해 이 규칙을 확인해 주었다. 릴레이 기계에는 그런 도구가 없고, 부품의
지연은 나노초가 아니라 밀리초이며, 같은 품번이라도 개체마다 지연이 두 배씩 다르다. 이 장은 그런 부품으로
틀림없이 동작하는 순서를 짜는 방법을 다룬다. 테스트 레지스터의 판정 순서, 그 순서를 지키게 하는 시간
규칙, 그리고 두 기계가 그 규칙을 지키는 방법(릴레이식의 순서기, 기어식의 슬롯 원판)이다.

\section{들쭉날쭉한 부품으로 시간을 정하기}\label{sec:t-bounds}

\Why 이 프로젝트가 쓰는 릴레이 시간은 표~\ref{tab:t-relay}와 같다. 데이터시트에 적힌 것은 대부분 ``최대''이고,
최소는 알 수 없는 경우가 많다. 그러면 시간 규칙은 모두 \emph{하한}으로 적어야 한다. ``\code{STROBE}는 7.5~ms
\emph{이상}'', ``\code{JUDGE}는 \code{STROBE}가 오른 뒤 11.5~ms \emph{이후}''처럼. 하한으로 적은 규칙에는 좋은
성질이 있다. 기계가 \emph{느려지면} 언제나 더 잘 지켜진다. 손으로 천천히 돌리든, 관람객이 한 걸음씩
누르든, 모터가 가속하는 중이든 상관없다. 위험한 것은 빨라지는 쪽 하나뿐이다. 이 장의 설계 결정은 대부분 이
성질에서 나온다.

\begin{table}[ht]
\centering
\caption{이 프로젝트가 쓰는 릴레이 시간(결정 0001). 일부는 데이터시트의 대표·최대값이고 일부는 추정이다.
시제품에서 실측해 확정한다.}
\label{tab:t-relay}
\small
\begin{tabular}{lrr}
\toprule
& 대표 & 범위 \\
\midrule
신호 릴레이가 붙는 시간(접점 안정 포함) & 2.5~ms & 2.5--5~ms \\
놓이는 시간, 코일에 다이오드만 & 3--5~ms & 3--10~ms (추정) \\
놓이는 시간, 다이오드 + 제너 30~V & 1.5~ms & 1.5--4~ms \\
접점이 공중에 떠 있는 틈(먼저 끊고 나중에 잇기) & & 0.2--0.6~ms (추정) \\
\bottomrule
\end{tabular}
\end{table}

\General 신호등의 노란불은 ``최소 3초''로 정한다. 교차로가 넓으면 늘이고, 짧게 하는 것만 금지한다. 엘리베이터의
문이 ``적어도 몇 초는 열려 있어야 한다''는 것도 같다. 둘 다 사람이 얼마나 빨리 움직일지 정확히 모르므로
가장 느린 사람에 맞춘 하한을 둔다.

\Ours 기어식은 빠른 드럼 60~rpm에서 한 위치가 38.46~ms이고, 릴레이식은 순서기가 한 위치를 63~ms로 낸다.
두 기계 모두 같은 테스트 레지스터(공통부)를 쓰므로, 레지스터가 요구하는 하한은 두 기계가 함께 지킨다.

\section{테스트 레지스터: 붙잡아 두고 나중에 판정한다}\label{sec:t-register}

테스트 레지스터는 봄베의 판단 장치다. \ref{sec:bombe-one-page}절에서 본 대로, 테스트 글자 케이블의 26선이
\emph{전부} 통전되면 그 위치는 버리고, 하나라도 통전되지 않으면 정지한다. 26선마다 \term{감지 릴레이}{sense
relay} $K$가 하나씩 있다.

\Why 처음 회로는 판정을 시험 전압이 걸려 있는 동안 끝내려 했다. 시험 전압이 걸리면 감지 릴레이가 붙고, 26개가
모두 붙으면 직렬로 이은 접점 사슬이 닫혀 \term{통과 릴레이}{pass relay} $P$가 붙고, 판정 순간에 $P$가 놓여
있으면 \term{정지 릴레이}{stop relay} $S$가 붙는다. 릴레이 세 단이 차례로 움직여야 한다. 최악의 시간으로 더하면
$5 + 5 + 5 = 15$~ms인데, 기어식에서 드럼 접점이 패드 위에 확실히 얹혀 있는 창은 6~ms다. 대표값으로도
8~ms라 들어가지 않는다. 기어를 늦추면 들어가지만, 최악값을 담으려면 속도가 3분의 1이 된다.

해법은 \term{표본 유지}{sample-and-hold}다. 시험 전압이 걸린 짧은 창 안에서는 감지 릴레이가 \emph{붙기만} 하고,
붙은 릴레이는 자기 유지(\ref{sec:relay-latch}절)로 그 결과를 붙든다. 판정은 창이 닫힌 \emph{뒤에} 느긋하게 한다.
창 안에서 끝나야 하는 것은 릴레이 하나의 동작(최악 5~ms)뿐이다. 판정 순간에 감지 릴레이들은 이미 멈춰 있는
값이므로 빠르기 경쟁도 없다.

\General 디지털카메라는 셔터가 열린 짧은 순간에 빛을 받아 두고, 사진 처리는 셔터가 닫힌 뒤에 한다. 아날로그-디지털
변환기 앞의 표본 유지 회로도 같다. 변하는 입력을 한 순간 커패시터에 붙잡아 두고, 변환은 그 붙잡힌 값으로
천천히 한다.

\Ours 그래서 뱅크(기어식 구동부나 릴레이식 순서기)는 테스트 레지스터에 세 신호를 차례로 보낸다.

\begin{itemize}
\item \code{STROBE}: 시험 전압을 메뉴 배선망에 건다. 통전된 선의 $K$가 붙어 스스로를 유지한다.
\item \code{JUDGE}: 판정 경로를 연다. 그때 $P$가 놓여 있으면(26선이 전부는 아니면) $S$가 붙어 정지를 붙든다.
\item \code{CLEAR}: 유지 버스 $H$를 끊어 $K$와 $P$를 놓는다. 다음 위치를 위한 원상 복귀다.
\end{itemize}

세 신호는 ``언제''만 알린다. ``정지인가''는 레지스터 안의 릴레이 접점이 정한다. 그림~\ref{fig:t-cycle}은 한
위치 동안 일어나는 일이다.

\begin{figure}[ht]
\centering
\begin{tikztimingtable}[timing/slope=0.05,timing/coldist=4pt,xscale=0.36,yscale=1.15,semithick]
  \code{STROBE} (접점) & 6H 32.5L \\
  $K$ (통전된 선) & 2.5L 28.5H 7.5L \\
  $P$ (26선 전부일 때) & 7L 23H 8.5L \\
  \code{JUDGE} (접점) & 12L 7H 19.5L \\
  $S$ (정지일 때) & 16L 22.5H \\
  \code{CLEAR} (접점) & 21L 12H 5.5L \\
\extracode
  \begin{pgfonlayer}{background}
    \foreach \x in {0,6,12,19,21,33,38.5} \draw[cgray!50,dashed] (\x,1.2) -- (\x,-11.5);
  \end{pgfonlayer}
  \foreach \x in {0,6,12,19,21,33,38.5} \node[font=\tiny,text=cgray] at (\x,-12.0) {\x};
  \node[note,anchor=west] at (39.2,-12.0) {ms};
  \node[note] at (3,1.7) {감지};
  \node[note] at (9,1.7) {정착};
  \node[note] at (15.5,1.7) {판정};
  \node[note] at (27,1.7) {복귀};
\end{tikztimingtable}
\caption{테스트 레지스터의 한 위치(기어식 60~rpm, 38.5~ms, 접점이 실제로 닫힌 구간). $P$와 $S$는 둘 가운데
하나만 일어난다. 26선이 전부 통전되면 $P$가 붙어 판정 경로를 막고, 아니면 $P$가 놓인 채라 $S$가 붙는다.
$S$가 붙으면 \code{CLEAR}가 와도 $K$가 지워지지 않는다. 그림의 $K$는 정지하지 않은 경우다.}
\label{fig:t-cycle}
\end{figure}

그림~\ref{fig:t-sh}는 이 동작의 핵심만 추린 회로다. 세 가지를 눈여겨볼 만하다.

\begin{figure}[ht]
\centering
\begin{circuitikz}[x=1cm,y=1cm]
  \ctikzset{bipoles/length=9mm}
  \draw (0,3) node[left,font=\footnotesize] {$+24$ V} to[nos,l={\scriptsize \code{STROBE}}] (1.6,3) -- (2.0,3);
  \node[cloud,minimum width=26mm,minimum height=8mm,anchor=west] (net) at (2.0,3) {메뉴 배선망};
  \draw (net.east) -- (5.6,3) node[above,font=\footnotesize] {선 $L_A$} -- (6.4,3);
  \draw (6.4,3) to[D] (6.4,1.6) coordinate (n);
  \node[rcoil,anchor=north] (K) at (6.4,1.3) {$K_A$};
  \draw (n) -- (K.north); \draw (K.south) -- (6.4,0.0) node[below,font=\footnotesize] {0 V};
  \draw (0,1.6) node[left,font=\footnotesize] {$+24$ V} to[ncs,l={\scriptsize \code{CLEAR}의 $r$}] (2.2,1.6)
        -- (3.0,1.6) node[above,font=\footnotesize,text=cgreen] {$H$} -- (3.8,1.6) to[nos,l={\scriptsize $k_A$ (자기 유지)}] (5.6,1.6) -- (n);
  \node[circ] at (n) {};
  \node[note,anchor=west,align=left] at (7.3,2.4) {차단 다이오드: 유지 전압이\\배선망으로 거꾸로 가지 않는다};
  \node[note,anchor=west,align=left] at (7.3,0.6) {$K_A$가 붙으면 $k_A$가 닫혀\\$H$에서 스스로를 먹인다};
\end{circuitikz}
\caption{감지 릴레이 하나의 표본 유지(26개 가운데 하나). 실제 회로에는 이 밖에 접점 사슬, 통과 릴레이 $P$,
정지 릴레이 $S$, 해제 릴레이가 있다(그림~\ref{fig:t-register-full}).}
\label{fig:t-sh}
\end{figure}

\begin{itemize}
\item \textbf{유지 버스 $H$.} 26개 감지 릴레이의 자기 유지 접점이 모두 같은 버스 $H$에서 전원을 받는다. $H$의
공급은 \code{CLEAR} 릴레이의 NC 접점을 지난다. \code{CLEAR} 동안 그 접점이 열리면 26개가 한꺼번에 놓인다.
정지 중에는 $S$의 접점이 이 NC 접점을 우회해 $H$를 살려 두므로, 정지 순간의 26선 상태가 램프에 남는다.
\item \textbf{차단 다이오드.} 감지 릴레이가 자기 유지에 들면 코일 쪽 마디에 +24~V가 계속 걸린다. 이 마디가
선에 그대로 이어져 있으면 그 전압이 메뉴 배선망으로 거꾸로 들어가 다른 선을 통전시키고, \code{STROBE}가
끝난 뒤에도 드럼 접점에 전류가 흐른다. 선마다 다이오드 하나가 이 길을 막는다. \ref{sec:relay-logic}절의 샛길
문제를 다이오드로 푼 예다.
\item \textbf{세 신호의 역할 분담.} \code{STROBE}, \code{JUDGE}, \code{CLEAR}는 레지스터 안의 타이밍 릴레이 셋을
움직인다. 시험 전류를 잇고 끊는 것도, 판정 경로를 여는 것도, $H$를 끊는 것도 레지스터 안의 접점이다.
\end{itemize}

\begin{lesson}{계산으로 정한 회로를 시뮬레이션이 고쳤다}
결정 0001은 처음에 계산으로만 이 회로를 정했다. 보드를 회로도로 옮기고 넷리스트를 시간 있는 스위치 수준
시뮬레이션(\ref{ch:boardverify}장)으로 돌리자 틀린 곳 넷이 나왔다. 그 가운데 둘은 \emph{놓친 정지}(위험한 쪽)로
이어지는 것이었다. 하나는 접점 사슬을 +24~V에서 먹인 것이다. 26선이 전부 붙은 위치에서는 \code{CLEAR} 동안에도
$P$의 코일이 사슬로 전압을 받아, $P$가 첫 $K$가 놓인 뒤에야 놓이기 시작했다. 그래서 다음 위치까지 $P$가 남아
그 위치를 통과로 만들었다. 사슬을 $H$에서 먹이게 고치자 $K$와 $P$가 같은 순간에 놓이기 시작했다. 다른 하나는
해제 버튼이 $S$만 풀던 것이다. 정지 중에 그 위치의 \code{CLEAR}가 이미 지나갔으면, 재출발 뒤 첫 \code{STROBE}가
지워지지 않은 $K$ 위에 덧붙어 다음 위치의 정지를 놓친다. 해제 릴레이 하나를 더해 해제 때 $H$까지 끊게 했다.
\end{lesson}

보드의 실제 회로는 그림~\ref{fig:t-register-full}이다(보고서와 같은 그림). 시뮬레이션은 릴레이가 가장 느릴 때
레지스터가 견디는 최소 폭도 찾았다. 접점 면에서 \code{STROBE} 5~ms(감지 릴레이의 최악 동작), \code{JUDGE}
5~ms(정지 릴레이의 최악 동작, 모자라면 정지를 놓친다), \code{CLEAR} 10~ms(감지 릴레이의 최악 복귀, 모자라면
$H$가 살아날 때 $K$가 다시 붙어 지워지지 않는다)다.

\begin{figure}[p]
\centering
\resizebox{\linewidth}{!}{% 테스트 레지스터와 정지 릴레이의 회로도 (06-hardware.tex, 그림 fig:hw-register)
% 코일은 대문자(K_A, P, S), 그 릴레이의 접점은 소문자(k_A, p, s)로 적는다.
\begin{circuitikz}[x=1cm,y=1cm]
  \ctikzset{bipoles/length=9mm}
  % ---- (1) 시험 전류: +24 V, STROBE 접점, 메뉴 배선망, 감지 릴레이 코일 26개 ----
  \node[font=\footnotesize,anchor=east] (v1) at (0,6.6) {$+24$ V};
  \draw (v1.east) to[nos,l={\footnotesize \code{STROBE}}] (2.3,6.6) -- (2.8,6.6);
  \node[cloud,minimum width=68mm,minimum height=9mm,anchor=west] (net) at (2.8,6.6)
    {메뉴 배선망 (스크램블러, 대각선 보드)};
  \foreach \x/\l in {3.5/A,5.0/B,6.5/C,8.9/Z} {
    \node[rcoil] (K\l) at (\x,5.2) {$K_{\l}$};
    \draw (\x,6.15) -- (K\l.north);
    \draw (K\l.south) -- (\x,4.3);
  }
  \node at (7.7,5.2) {$\cdots$};
  \draw (2.9,4.3) -- (9.6,4.3) node[right,font=\footnotesize] {0 V};
  \node[note,anchor=west] at (10.3,6.0) {입력 케이블의 26선마다\\감지 릴레이 코일 하나};
  \node[font=\footnotesize\bfseries,anchor=west] at (-1.1,7.4) {(1) 시험 전류의 경로};
  % ---- (2) 접점 사슬과 통과 릴레이 ----
  \node[font=\footnotesize\bfseries,anchor=west] at (-1.1,3.5) {(2) 접점 사슬: 26선이 전부 통전되었는가};
  \node[font=\footnotesize,anchor=east] (v2) at (0,2.6) {$+24$ V};
  \draw (v2.east) -- (2.6,2.6) to[nos,l={$k_A$}] (4.1,2.6) to[nos,l={$k_B$}] (5.6,2.6)
        to[nos,l={$k_C$}] (7.1,2.6);
  \draw[dashed] (7.1,2.6) -- (8.1,2.6);
  \draw (8.1,2.6) to[nos,l={$k_Z$}] (9.6,2.6) -- (10.6,2.6);
  \node[rcoil,anchor=west] (P) at (10.6,2.6) {$P$};
  \draw (P.east) -- (12.3,2.6) node[right,font=\footnotesize] {0 V};
  \node[note,anchor=north] at (11.05,2.25) {통과 릴레이};
  % ---- (3) 정지 릴레이와 자기 유지 ----
  \node[font=\footnotesize\bfseries,anchor=west] at (-1.1,1.55) {(3) 정지 릴레이};
  \node[font=\footnotesize,anchor=east] (v3) at (0,0.7) {$+24$ V};
  \draw (v3.east) -- (0.8,0.7) coordinate (a) -- (2.6,0.7) to[nos,l={$j$}] (4.1,0.7)
        to[ncs,l={$p$}] (5.6,0.7) -- (8.6,0.7) coordinate (b) -- (10.6,0.7);
  \node[rcoil,anchor=west] (S) at (10.6,0.7) {$S$};
  \draw (S.east) -- (12.3,0.7) node[right,font=\footnotesize] {0 V};
  \node[note,anchor=north] at (11.05,0.35) {정지 릴레이};
  \draw (a) -- (0.8,-0.7) -- (2.6,-0.7) to[ncs,l_={원위치 버튼}] (4.1,-0.7) -- (5.9,-0.7)
        to[nos,l_={$s$ (자기 유지)}] (7.4,-0.7) -- (8.6,-0.7) -- (b);
  \node[circ] at (a) {};
  \node[circ] at (b) {};
  % ---- (4) RUN ----
  \node[font=\footnotesize\bfseries,anchor=west] at (-1.1,-1.75) {(4) 구동 허가};
  \node[font=\footnotesize,anchor=east] (v4) at (0,-2.5) {$+24$ V};
  \draw (v4.east) -- (2.6,-2.5) to[nos,l_={시작}] (4.1,-2.5) to[ncs,l_={$s$}] (5.6,-2.5) -- (6.6,-2.5)
        node[right,font=\footnotesize] {뱅크의 \code{RUN} (정지 릴레이가 동작하면 끊긴다)};
\end{circuitikz}
}
\caption{테스트 레지스터 보드의 회로(보고서의 그림과 같다). 코일은 대문자($K_A$, $P$, $S$), 그 릴레이의
접점은 소문자다. $t$, $j$, $r$은 \code{STROBE}, \code{JUDGE}, \code{CLEAR}가 움직이는 타이밍 릴레이의 접점,
$x$는 해제 릴레이의 접점이다.}
\label{fig:t-register-full}
\end{figure}

\section{두 면: 커넥터의 전압과 접점의 실제}\label{sec:t-faces}

\Why 레지스터가 요구하는 시간은 \emph{레지스터 안의 접점}이 닫혀 있는 구간에 대한 것이다. 그런데 뱅크가 낼 수
있는 것은 \emph{커넥터 핀의 전압}뿐이다. 둘 사이에 타이밍 릴레이가 하나 있고, 그 릴레이는 핀이 오른 뒤
2.5--4~ms에 붙고(켜짐 지연 $d_{\mathrm{on}}$), 핀이 내린 뒤 1.5--4~ms에 놓인다(꺼짐 지연 $d_{\mathrm{off}}$).
결정 0003은 이 둘을 \term{접점 면}{contact face}과 \term{커넥터 면}{connector face}이라는 이름으로 나눠, 회로의
조건은 접점 면에서 적고 뱅크가 지킬 규칙은 커넥터 면에서 적었다(그림~\ref{fig:t-faces}).

\begin{figure}[ht]
\centering
\begin{tikzpicture}[x=0.5cm,y=1cm]
  \draw[->] (-0.5,0) -- (22,0) node[right,font=\footnotesize] {시간};
  % connector face pulse
  \draw[line width=1pt,cblue] (0,1.6) -- (2,1.6) -- (2,2.4) -- (12,2.4) -- (12,1.6) -- (20,1.6);
  \node[font=\footnotesize,anchor=east,text=cblue] at (-0.6,2.0) {커넥터 면 (핀)};
  % contact face pulse (range)
  \fill[cbluebg] (6,0.4) rectangle (10,1.0);
  \fill[cyellowbg] (12,0.4) rectangle (16,1.0);
  \draw[line width=1pt,corange] (0,0.4) -- (6,0.4) -- (6,1.0) -- (14,1.0) -- (14,0.4) -- (20,0.4);
  \node[font=\footnotesize,anchor=east,text=corange] at (-0.6,0.7) {접점 면 (레지스터 안)};
  \draw[<->] (2,1.3) -- node[above,font=\scriptsize] {$d_{\mathrm{on}}$ 2.5--4} (6,1.3);
  \draw[<->] (12,1.3) -- node[above,font=\scriptsize] {$d_{\mathrm{off}}$ 1.5--4} (14,1.3);
  \node[note,anchor=north,align=center] at (8,0.3) {파란 띠: 접점이 붙는\\때가 흩어지는 범위};
  \node[note,anchor=north,align=center] at (14,0.3) {노란 띠: 놓이는\\때의 범위};
\end{tikzpicture}
\caption{두 면. 커넥터 핀의 펄스(위)는 타이밍 릴레이를 거쳐 접점의 펄스(아래)가 된다. 접점 면의 폭은 커넥터
면의 폭 $- d_{\mathrm{on}} + d_{\mathrm{off}}$이므로 최악의 경우($d_{\mathrm{on}}=4$, $d_{\mathrm{off}}=1.5$) 2.5~ms 줄어든다.}
\label{fig:t-faces}
\end{figure}

그러면 커넥터 면의 규칙은 접점 면의 조건에서 \emph{계산으로} 나온다. 예를 들어 보자.

\begin{itemize}
\item 접점 면에서 \code{STROBE}는 5~ms 이상 닫혀야 한다(C1). 폭은 최악 2.5~ms 줄므로, 핀에서는 $5 + 2.5 = 7.5$~ms
이상이어야 한다(L1). \code{JUDGE}(L3)도 같은 이유로 7.5, \code{CLEAR}(L5)는 $10 + 2.5 = 12.5$다.
\item 접점 면에서 \code{JUDGE}는 \code{STROBE}가 닫힌 뒤 10~ms(감지 5 + 통과 5) 이후에 닫혀야 한다(C2).
\code{STROBE}의 접점은 핀이 오른 뒤 \emph{늦으면} 4~ms에 닫히고, \code{JUDGE}의 접점은 핀이 오른 뒤
\emph{빠르면} 2.5~ms에 닫힌다. 가장 나쁜 짝을 맞추면 \code{JUDGE}의 핀은 \code{STROBE}의 핀보다 $4 + 10 - 2.5 =
11.5$~ms 뒤에 올라야 한다(L2).
\item \code{CLEAR}의 접점은 \code{JUDGE}의 접점이 열린 뒤에 열려야 한다(C4). \code{JUDGE}의 접점은 핀이 내린 뒤
늦으면 4~ms에 열리고, \code{CLEAR}의 접점은 핀이 오른 뒤 빠르면 2.5~ms에 열린다. 그래서 $4 - 2.5 = 1.5$~ms(L4).
\end{itemize}

이렇게 얻은 규칙이 표~\ref{tab:t-rules}의 L1--L9다. 모두 하한이거나 순서이고, 레지스터의 실제 지연을 몰라도
성립한다. 하한을 모두 더한 가장 짧은 주기는 $0 \to 11.5 \to 19 \to 20.5 \to 33 \to 34.5$, 곧 34.5~ms다.

\begin{table}[ht]
\centering
\caption{뱅크가 지킬 규칙(결정 0003, 커넥터 면, ms).}
\label{tab:t-rules}
\small
\begin{tabularx}{\linewidth}{>{\raggedright\arraybackslash}p{9mm}>{\raggedright\arraybackslash}X>{\raggedright\arraybackslash}p{38mm}}
\toprule
& 요구 & 어기면 \\
\midrule
L1 & \code{STROBE}의 폭 $\ge 7.5$ & 감지 릴레이가 못 붙음(거짓 정지) \\
L2 & \code{JUDGE}의 오름 $\ge$ \code{STROBE}의 오름 $+ 11.5$ & $P$가 붙기 전에 판정(거짓 정지) \\
L3 & \code{JUDGE}의 폭 $\ge 7.5$ & $S$가 못 붙음(\textbf{놓친 정지}) \\
L4 & \code{CLEAR}의 오름 $\ge$ \code{JUDGE}의 내림 $+ 1.5$ & 판정 중에 $P$가 놓임 \\
L5 & \code{CLEAR}의 폭 $\ge 12.5$ & $K$가 지워지지 않음(\textbf{놓친 정지}) \\
L6 & 다음 \code{STROBE}의 오름 $\ge$ \code{CLEAR}의 내림 $+ 1.5$ & 유지 버스가 없는 동안 붙은 $K$를 잃음 \\
L7 & \code{STROBE}를 낸 위치에는 \code{JUDGE}와 \code{CLEAR}를 반드시 낸다 & 판정 없이 지나간 위치 \\
L8 & 전원 투입, \code{HOME}, 해제 뒤에는 \code{STROBE} 전에 \code{JUDGE}를 내지 않는다 & 빈 레지스터를 판정 = 정지 \\
L9 & 시프터·드럼의 전환은 \code{CLEAR}의 오름부터 다음 \code{STROBE}의 오름 사이에만 & 전류가 흐르는 접점을 옮김 \\
\bottomrule
\end{tabularx}
\end{table}

L9는 \term{무전류 전환}{dry switching} 규칙이다. 전류가 흐르는 채로 접점을 떼면 불꽃이 튀어 접점이 상한다.
시험 전압이 없는 동안에만 드럼이나 시프터를 움직이라는 것이다.

\General 두 면의 구분은 디지털 회로의 \term{셋업 시간}{setup time}과 같은 생각이다. 칩의 데이터시트는 ``핀에서
클럭 에지 몇 ns 전에 데이터가 안정되어야 한다''고 적는다. 칩 안의 실제 플립플롭이 요구하는 조건을, 핀에서 그
플립플롭까지의 지연이 가장 나쁠 때를 넣어 핀의 조건으로 옮긴 것이다.

\Ours 결정 0003은 하나 더 열어 두었다. 물린 레지스터의 지연을 실제로 재서(\term{보정}{calibration}) 그만큼
앞당겨 내면 접점 면의 조건을 직접 지킬 수 있다. 기어식은 \code{STROBE}만 이렇게 보정하고(드럼 접점의 가운데에
접점 면의 창을 놓아야 하므로), 나머지는 보정 없이 넓게 낸다. 릴레이식 순서기는 셋 모두 보정 없이 지킨다.

\section{릴레이식: 555와 링 계수기가 박자를 친다}\label{sec:t-sequencer}

릴레이식 뱅크의 \term{순서기}{sequencer} 보드는 위치마다 \code{STROBE},
\code{JUDGE}, \code{CLEAR}, 그리고 계수기의 \code{CK}/\code{NCK}를 L1--L9에 맞춰 낸다. 정지와 재출발,
\code{HOME}(계수기 프리셋), 마지막 위치(\code{END}), 대기도 이 보드가 처리한다. 릴레이 50개가 있다(폴더
\code{hardware/relay/sequencer}).

\Why 순서·조건·잠금은 모두 릴레이 접점으로 만들었다. 그런데 \emph{박자}, 곧 ``6.3~ms마다 다음 단계''라는 시간
기준 하나만은 전자식이다. 이유는 \ref{sec:t-bounds}절의 하한에 있다. 규칙이 하한이면 시간 기준의 \emph{최소}를
보증할 수 있어야 한다. 릴레이로 지연을 만들려면 코일에 커패시터를 달아 천천히 충전되게 하는데, 릴레이가 붙는
전압은 ``정격의 80\% 이하''처럼 위쪽 한계만 알려져 있다. 실제로 50\%에서 붙는 개체와 80\%에서 붙는 개체는 같은
RC에서 지연이 두 배 넘게 다르다. 하한을 보증할 수 없다. 릴레이를 사슬로 이어 박자를 만들면 단마다 흩어짐이
쌓여 한 주기가 70--80~ms로 늘어난다. 그래서 전자 부품 두 개로 박자를 만들고, 그 박자가 정하는 것은 ``언제''뿐이며
``진행할지''는 여전히 릴레이가 정한다. 기어식에서 슬롯 원판이 하는 일과 같다(\ref{sec:t-gear}절).

\General 시간 기준은 입문 전자 회로의 단골 둘이다.
\begin{itemize}
\item \textbf{555 타이머}는 커패시터 하나를 저항을 거쳐 충전하다가 전원의 2/3에 이르면 방전으로 바꾸고, 1/3까지
내려가면 다시 충전으로 바꾸는 칩이다. 출력은 충전 중에 높고 방전 중에 낮은 네모파가 된다. 한 주기는
\[
  T = 0.693\,(R_1 + 2R_2)\,C
\]
이다. 충전은 $R_1 + R_2$를 지나고 방전은 $R_2$만 지나기 때문에 $R_2$가 두 번 들어간다.
\item \textbf{74HC4017}은 클럭이 올 때마다 출력 Q0, Q1, \ldots, Q9 가운데 \emph{하나만} 차례로 켜는 칩이다(10단
링 계수기). 흐르는 불빛 장식을 만드는 입문 예제에 꼭 나온다.
\end{itemize}

\Ours 순서기는 LMC555($R_1$ = 10~k$\Omega$, $R_2$ = 150~k$\Omega$ + 트리머 100~k$\Omega$, $C$ = 22~nF)가 6.3~ms마다
4017을 한 칸 민다. 열 칸이 한 위치이므로 한 주기는 63~ms다. 4017의 출력 몇 개를 다이오드 OR로 묶어
\term{창}{window}을 만들고, 창이 고측 구동기(\ref{sec:el-mosfet}절)를 거쳐 출력 릴레이의 코일을 먹인다(그림~\ref{fig:t-ring}).

\begin{figure}[ht]
\centering
\begin{tikzpicture}[x=1.15cm,y=0.62cm]
  \foreach \q in {0,...,9} {
    \draw[cgray!60] (\q,0.6) -- (\q,-5.6);
    \node[font=\scriptsize\ttfamily] at (\q+0.5,1.0) {Q\q};
  }
  \draw[cgray!60] (10,0.6) -- (10,-5.6);
  \newcommand{\seqwin}[5]{\fill[#4] (#1,#3-0.35) rectangle (#2,#3+0.35); \node[font=\scriptsize,anchor=east] at (-0.15,#3) {#5};}
  \seqwin{0}{2}{0}{cbluebg}{\code{STROBE} (Q0--Q1)}
  \seqwin{3}{5}{-1.2}{cgreenbg}{\code{JUDGE} (Q3--Q4)}
  \seqwin{6}{9}{-2.4}{corangebg}{\code{CLEAR} (Q6--Q8)}
  \seqwin{6}{7}{-3.6}{cyellowbg}{\code{CK} 결정 (Q6)}
  \seqwin{7}{9}{-4.8}{credbg}{\code{CK} (Q7--Q8)}
  \node[font=\scriptsize,anchor=west] at (9.05,-4.8) {Q9: 주차};
  \draw[<->] (0,-6.2) -- node[below,font=\scriptsize] {한 틱 $T$ = 6.3 ms (관찰 모드 27 ms)} (1,-6.2);
  \draw[<->] (0,-7.4) -- node[below,font=\scriptsize] {한 위치 = 10틱 = 63 ms (관찰 모드 270 ms)} (10,-7.4);
\end{tikzpicture}
\caption{순서기의 링 상태와 창. 창이 바로 출력이 되지는 않는다. 창과 출력 릴레이의 코일 사이에 상태 릴레이의
접점(순환 허가, \code{CK} 결정 등)이 직렬로 들어가, 접점이 닫혀 있어야만 창이 코일에 닿는다.}
\label{fig:t-ring}
\end{figure}

왜 63~ms인가. 출력이 릴레이 접점이므로 순서기 \emph{자신의} 출력 릴레이도 붙고 놓이는 데 시간이 든다.
\code{STROBE}의 핀 폭은 구동 창의 폭에서 (붙는 시간 $-$ 놓이는 시간)만큼 깎여 최악 3.5--4.1~ms 줄어든다. 그래서
L1의 7.5~ms를 지키려면 구동 창이 11.6~ms 이상, 곧 두 틱이 5.8~ms 이상이어야 하고, 여유를 두어 6.3~ms로 정했다.
시뮬레이션이 릴레이 시간을 가장 나쁘게 넣어 찾은 \code{STROBE}의 최소 폭은 8.23~ms(여유 0.73~ms)였다.

\code{CK}를 낼지는 Q6에서 정한다. 그때 \code{RUN}(정지가 아님)이 높으면 \code{CK} 결정 릴레이가 붙어 Q6--Q8 동안
스스로를 붙들고, Q7--Q8에 \code{CK}가 나간다. 정지가 나서 \code{RUN}이 내려가면, 릴레이 시간이 가장 나쁜
조합에서도 Q6이 시작된 뒤 0.6~ms 안에 \code{RUN}을 따르는 릴레이가 열려, 붙는 데 2.5~ms 이상 걸리는 결정
릴레이의 동작이 취소된다(시뮬레이션으로 확인). 이 순서가 중요하다. 계수기는 내려 셀 수 없으므로 정지한 위치를 한 번 지나치면 되돌아올
방법이 없다.

\subsection{관찰 모드: 왜 커패시터가 아니라 저항을 바꾸는가}

전시에서는 관람객이 릴레이가 넘어가는 것을 눈으로 따라갈 수 있게 느린 \term{관찰 모드}{observation mode}가
필요하다. 순서기는 555의 $R_2$에 680~k$\Omega$을 직렬로 넣어 틱을 27.0~ms로 늘린다(초당 3.7 위치). 평소에는
릴레이 접점 두 개가 직렬로 이 저항을 단락해 두고, 관찰 쪽으로 스위치를 돌리면 단락이 열린다.

\Why 속도를 바꾸는 순간은 운전 중 아무 때나다. 관람객이 손잡이를 돌리기 때문이다. 하한 규칙에서 위험한 것은
틱이 \emph{짧아지는} 것 하나뿐이므로, 바꾸는 순간이 걸린 틱도 빠른 틱보다 짧아지면 안 된다(그림~\ref{fig:t-555}).

\begin{itemize}
\item \textbf{저항을 바꾸면} 커패시터의 전압은 그대로 이어지고 충전·방전의 빠르기만 느려진다. 전압이 빠른
모드보다 빨리 움직일 수 없으므로 어느 틱도 빠른 틱보다 짧아지지 않는다. 링은 한 번 Q9를 떠나면 반드시 Q9까지
돌므로 주기가 중간에 끊기지도 않는다.
\item \textbf{커패시터를 더하면} 충전된 커패시터와 빈 커패시터가 전하를 나눠 전압이 순간 떨어진다. 방전 중이었다면
곧바로 1/3 아래로 내려가 555가 충전으로 넘어간다. 그 틱이 반쯤에서 잘려 \code{STROBE}나 \code{JUDGE}의 폭이
규칙보다 짧아질 수 있다.
\item \textbf{분주기}(4017을 하나 더 두어 느린 클럭을 고르는 것)도 두 클럭을 비동기로 바꾸는 순간 틱이 임의로
짧아진다. 막으려면 동기화 회로가 필요해 전자 부품이 는다.
\end{itemize}

\begin{figure}[ht]
\centering
\begin{tikzpicture}[x=1cm,y=2.2cm]
  \foreach \dx/\ttl in {0/저항을 바꿈,7.6/커패시터를 더함} {
    \begin{scope}[xshift=\dx cm]
      \node[font=\small\bfseries] at (3,1.25) {\ttl};
      \draw[->] (0,0) -- (6.4,0) node[right,font=\scriptsize] {$t$};
      \draw[->] (0,0) -- (0,1.1) node[above,font=\scriptsize] {$v_C$};
      \draw[cgray,dashed] (0,1/3) -- (6.2,1/3) node[right,font=\tiny] {1/3};
      \draw[cgray,dashed] (0,2/3) -- (6.2,2/3) node[right,font=\tiny] {2/3};
      % two normal ticks: charge (tau 1) 1/3->2/3 takes ln2 = 0.693, discharge same
      \draw[cblue,line width=0.9pt,domain=0:0.693,samples=20] plot (\x,{1-(2/3)*exp(-\x)});
      \draw[cblue,line width=0.9pt,domain=0:0.693,samples=20] plot (\x+0.693,{(2/3)*exp(-\x)});
      \draw[cblue,line width=0.9pt,domain=0:0.693,samples=20] plot (\x+1.386,{1-(2/3)*exp(-\x)});
      \draw[cred] (2.379,-0.05) -- (2.379,1.05);
      \node[font=\tiny,text=cred,anchor=south] at (2.379,1.05) {바꾸는 순간};
    \end{scope}}
  % left: discharge starts at 2.079 with tau 1 for 0.3, then tau 4
  \draw[cblue,line width=0.9pt,domain=0:0.3,samples=10] plot (\x+2.079,{(2/3)*exp(-\x)});
  \draw[corange,line width=0.9pt,domain=0:1.57,samples=30] plot (\x+2.379,{(2/3)*exp(-0.3)*exp(-\x/4)});
  \draw[corange,line width=0.9pt,domain=0:0.9,samples=20] plot (\x+3.949,{1-(2/3)*exp(-\x/4)});
  \node[note,anchor=west,align=left,text=corange] at (2.5,0.15) {전압은 이어지고\\내려가는 빠르기만 느려진다};
  % right: discharge 0.3 then voltage halves (equal empty capacitor added) -> below 1/3 -> charging restarts
  \begin{scope}[xshift=7.6cm]
    \draw[cblue,line width=0.9pt,domain=0:0.3,samples=10] plot (\x+2.079,{(2/3)*exp(-\x)});
    \draw[corange,line width=0.9pt] (2.379,{(2/3)*exp(-0.3)}) -- (2.379,{(1/3)*exp(-0.3)});
    \draw[corange,line width=0.9pt,domain=0:1.6,samples=30] plot (\x+2.379,{1-(1-(1/3)*exp(-0.3))*exp(-\x/2)});
    \node[note,anchor=west,align=left,text=cred] at (2.6,0.12) {1/3 아래로 떨어져\\틱이 반쯤에서 잘린다};
  \end{scope}
\end{tikzpicture}
\caption{555 커패시터 전압의 모양(개념도). 왼쪽은 방전 도중에 $R_2$를 늘린 경우로, 그 틱은 오히려 길어진다.
오른쪽은 같은 순간에 빈 커패시터를 병렬로 더한 경우로, 전압이 순간 반으로 떨어져 555가 곧바로 충전으로
넘어간다. 그 틱은 짧아진다.}
\label{fig:t-555}
\end{figure}

\Ours 순서기의 시뮬레이터는 555를 넷리스트의 RC로 흉내 낸다. 그 순간 닫혀 있는 접점으로 이어진 저항의 합을
$R_2$로 쓰고, 저항이 바뀌면 커패시터 전압을 이어서 계산한다(\ref{sec:v-time}절). 운전 중 선택 스위치를 무작위
시각에 36번씩 뒤집어도 가장 짧은 555 주기가 빠른 틱의 0.997배였다(넣어 둔 흔들림 $\pm$0.3\% 안). 반대로
``관찰 저항이 들어가지 않는'' 고장을 일부러 넣은 회로는 주기가 59.9~ms로 나와 검증기에 걸린다.

\subsection{한 걸음 버튼}

관찰 모드보다 더 느린 것이 \term{한 걸음}{single step} 버튼이다. 한 번 누르면 정확히 한 위치만 진행한다.

\Why ``정확히 한 바퀴''는 이미 4017이 보장한다. 4017은 Q9에서만 서도록(주차) 배선되어 있어, 링이 한 번 Q9를
떠나면 \code{STROBE}, \code{JUDGE}, \code{CLEAR}, \code{CK}를 모두 낸 뒤 Q9로 돌아와 선다. 그래서 버튼 회로가 할
일은 ``한 바퀴를 시작하게 하고, 시작했으면 다음 바퀴는 막는 것''뿐이다. 이것을 릴레이 두 개의 래치로 만든다.
\begin{itemize}
\item \textbf{받아들임} 릴레이(\code{KSA}): 링이 주차해 있을 때 버튼이 눌리면 선다. 링을 출발시키는 조건(GO)에
들어간다.
\item \textbf{주기} 릴레이(\code{KSC}): 받아들인 한 걸음의 링이 Q0에 들어서면 선다. 서자마자 그 NC 접점이 한
걸음의 GO를 끊는다. 이번 바퀴는 이미 출발했으므로 계속 돌고, 다음 바퀴는 출발하지 않는다.
\end{itemize}
버튼을 계속 누르고 있으면 두 릴레이가 서로를 붙들어 다음 바퀴가 없다(반복하지 않는다). 버튼을 일찍 놓아도
받아들임 릴레이는 링이 출발한 동안 다른 접점으로 붙들려 한 바퀴를 마친다.

\General 컴퓨터의 키보드나 마우스 버튼도 ``누를 때마다 정확히 한 번''이 필요하다. 접점은 닫힐 때 수 ms 동안
여러 번 튀기 때문에 그대로 세면 한 번 누름이 여러 번이 된다. 보통 펌웨어가 일정 시간 기다려 튐을 거른다.

\Ours 이 보드는 튐을 릴레이의 느림으로 거른다. 버튼을 따라가는 릴레이(\code{KSP})는 붙는 데 2.5~ms 이상이
걸리므로 그보다 짧게 닫히는 튐에는 붙지 않고, 코일에 다이오드만 달아 놓이는 데 3~ms 이상이 걸리므로 그보다
짧게 열리는 튐에는 놓이지 않는다. 받아들임은 주차 중에만 되므로 바퀴가 도는 동안의 튐은 아무 일도 하지 않는다.
한 걸음 주기는 선택 스위치와 무관하게 늘 관찰 속도로 돈다. 한 걸음 중에는 \code{RUN}이 낮으므로 정지가 나면
\code{STOP\_LATCH}를 거치는 조금 느린 길로 \code{CK}를 막아야 하는데, 빠른 틱에서는 그 여유가 0.1~ms뿐이기
때문이다.

\section{기어식: 슬롯 원판이 시계다}\label{sec:t-gear}

기어식에는 클럭이 따로 없다. 드럼이 도는 것 자체가 시간이다.

\Why 빠른 드럼과 기어로 묶여 도는 \term{주기 축}{cycle shaft}에 지름 160~mm의 얇은 원판이 있고, 원판에는 동심원 트랙 셋에
슬롯(구멍)이 뚫려 있다. 트랙마다 광센서 하나가 슬롯이 지나가는 동안 신호를 낸다. \code{STROBE}, \code{JUDGE},
\code{CLEAR}가 각각 한 트랙이다. 시각이 각도로 바뀐 것이어서, 속도가 느려지면 모든 구간이 시간으로 길어진다.
\ref{sec:t-bounds}절의 ``느려지면 늘 안전하다''가 기구로 성립하는 것이다.

주기 축 한 바퀴가 빠른 드럼 30칸이므로 한 위치가 주기 축의 12°이고, 60~rpm에서 1~ms가 0.312°다. 모터는 한 위치에
400스텝을 가므로 한 스텝이 0.096~ms다. 슬롯의 가장자리는 이 스텝의 정수배에 둔다. 26칸에만 슬롯이 있고 나머지
네 칸은 가운데 드럼이 한 칸 넘어가는 \term{자리올림}{carry} 구간이다.

\begin{figure}[ht]
\centering
\begin{tikzpicture}[x=0.29cm,y=1cm]
  \draw[->] (-4,0) -- (40,0) node[right,font=\scriptsize] {ms};
  \foreach \t in {0,10,20,30,38.46} \draw (\t,0.08) -- (\t,-0.08) node[below,font=\tiny] {\t};
  \fill[clight] (-3,0.25) rectangle (9,4.15);
  \node[note,anchor=south] at (3,4.15) {드럼 접점의 완전 접촉 ($-3$ … $+9$)};
  \newcommand{\slotbar}[5]{\fill[#4] (#1,#3-0.18) rectangle (#2,#3+0.18); \node[font=\scriptsize,anchor=east] at (-4.2,#3) {#5};}
  \slotbar{-2.50}{4.52}{3.6}{cbluebg}{\code{STROBE} 슬롯}
  \slotbar{0}{6}{3.15}{cblue!50}{접점 면}
  \slotbar{8.56}{16.63}{2.3}{cgreenbg}{\code{JUDGE} 슬롯}
  \slotbar{12}{19}{1.85}{cgreen!50}{접점 면}
  \slotbar{18.65}{31.73}{1.0}{corangebg}{\code{CLEAR} 슬롯}
  \slotbar{21}{33}{0.55}{corange!50}{접점 면}
  \fill[cbluebg] (35.96,3.42) rectangle (38.46,3.78);
  \draw[<->,cred] (31.73,1.45) -- node[above,font=\tiny,text=cred] {정지 자리} (35.96,1.45);
\end{tikzpicture}
\caption{기어식 한 위치의 시각표(60~rpm, \code{STROBE} 접점 시작이 0). 신호마다 위의 연한 띠가 슬롯(= 커넥터 면),
아래의 진한 띠가 결정 0001의 접점 면이다. 슬롯은 접점 면보다 앞서고 넓다. 타이밍 릴레이가 붙고 놓이는 시간을 미리 넣었기
때문이다. \code{CLEAR} 슬롯 끝과 다음 \code{STROBE} 슬롯 사이(44스텝, 4.23~ms)가 정지했을 때 축이 서는 자리다.}
\label{fig:t-slots}
\end{figure}

\Ours 슬롯의 각도는 L1--L9에서 계산했다. \code{JUDGE}와 \code{CLEAR}는 타이밍 릴레이의 지연이 범위 어디에 있어도
접점 면의 조건이 성립하도록 넓혔고, 각 규칙에 0.5~ms의 공차 여유(원판 가공, 센서 문턱, 전자 회로의 흩어짐을
더한 값)를 남겼다. 처음 판의 슬롯은 L4의 여유가 정확히 0이었다. 제어 보드의 시뮬레이션이 1.521~ms(하한
1.5)를 내놓자 기계 공차를 받을 수 없다는 것이 드러나 슬롯을 다시 잡았다. 슬롯 각도의 원본은 CAD의
\code{params.py} 한 곳이고, 제어 보드의 펌웨어 모델과 시뮬레이션이 같은 파일을 읽는다.

센서에서 커넥터 핀까지의 길에는 \emph{펌웨어가 없다}. 광센서 출력이 RC 필터와 슈미트 트리거(74HC14, \ref{sec:el-cap}절)를 지나
게이트 몇 개와 고측 구동기를 거쳐 핀으로 나간다. 마이크로컨트롤러는 모터의 가감속과 정지 위치로의 되돌림을
맡을 뿐, 이 길에는 ``지금 내보내도 된다''는 허가 신호 하나만 준다. 게이트는 두 규칙을 하드웨어로 지킨다.

\begin{lstlisting}[numbers=none,caption={기어식 제어 보드의 타이밍 경로(\code{control-board/README.md} 1.1절).},label={lst:t-gate}]
  STROBE_G = SLOT_S & TXE
  CLEAR_G  = SLOT_C & TXE
  JUDGE_G  = SLOT_J & SEEN & TXE
  TXE  = TXEN | CYC          TXEN = from the MCU (0 during HOME / moving back)
  CYC  : set while STROBE_G is high, cleared by the fall of CLEAR_G        (L7)
  SEEN : set by the rise of STROBE_G,
         cleared on power-up | HOMING | fall of STOP_LATCH                   (L8)
\end{lstlisting}

\code{CYC}는 ``\code{STROBE}를 냈으면 그 주기의 \code{CLEAR}까지는 낸다''(L7)를 지킨다. 펌웨어가 주기 한가운데서
허가를 거두어도 시작한 주기는 끝난다. \code{SEEN}은 ``\code{STROBE} 없이 \code{JUDGE}를 내지 않는다''(L8)를
지킨다. 비어 있는 레지스터에 \code{JUDGE}가 가면 곧 정지가 되기 때문이다.

두 기계를 나란히 놓으면 같은 생각이 두 번 나타난다. 시간 기준(슬롯 원판, 555와 4017)은 ``언제''만 정하고, 그
시간 기준이 아무리 흔들려도 틱이 짧아지지 않게 만들며, 순서와 조건은 각 기계의 방식대로(기어식은 게이트,
릴레이식은 릴레이 접점) 하드웨어가 지킨다.

\section{정리}

\begin{itemize}
\item 릴레이의 시간은 넓게 흩어지고 최소를 모른다. 그래서 시간 규칙은 모두 하한으로 적고, 기계가 느려지면 늘
더 잘 지켜지게 한다.
\item 테스트 레지스터는 표본 유지다. \code{STROBE} 동안 감지 릴레이가 붙어 스스로를 붙들고, \code{JUDGE}에서
판정하고, \code{CLEAR}에서 유지 버스를 끊어 지운다. 선마다 다이오드가 유지 전압의 역류를 막는다.
\item 회로의 조건은 접점 면에서, 뱅크의 규칙(L1--L9)은 커넥터 면에서 적는다. 둘 사이의 지연 범위를 가장 나쁘게
넣어 하나에서 다른 하나를 계산한다.
\item 릴레이식 순서기는 555와 4017로 박자만 만들고 나머지는 릴레이로 정한다. 관찰 모드는 커패시터 전압이 이어지는
저항 쪽을 바꿔 어떤 틱도 짧아지지 않게 한다. 한 걸음 버튼은 ``시작은 허락하고 다음 바퀴는 막는'' 래치다.
\item 기어식은 슬롯 원판이 시각을 각도로 갖고, 펌웨어 없는 게이트가 L7과 L8을 지킨다.
\end{itemize}
